\subsection*{68. Bayes Should Test Instructional Causality Separately From Demand}

A paid pilot can show that certain people were willing to purchase a defined experience and that the institution delivered something. It cannot, by itself, show that the teaching caused their entire improvement. A learning evaluation can show capability gain but not a profitable acquisition channel. A repeatable sales funnel can exist alongside weak educational outcomes.

Maintain at least 3 hypothesis ledgers: \textbf{commercial}, \textbf{learning}, and \textbf{operational}:

\paragraph*{The commercial ledger} asks which buyers pay, through what channel, and under what alternatives.

\paragraph*{The learning ledger} asks which capacities change and what evidence supports that inference.

\paragraph*{The operational ledger} asks whether the system can deliver repeatedly at viable economics and acceptable reviewer burden.

A founder who checks only 1 ledger may falsely conclude the institution has found its model.

\subsection*{69. One Possible Comparison Study After Initial Feasibility}

Once the tasks, rubric, service quality, and basic economics are stable, Bayes might compare 2 ethically acceptable ways to support the same learning objective: structured human review alone versus structured AI hints followed by human review; or synchronous worked examples versus prepared recorded examples with a live critique. If participants can be assigned fairly and both conditions provide responsible instruction, a randomized pilot might estimate a narrow, decision-relevant difference.

Predefine the primary independent assessment, the relevant cost per learner, balancing measures such as overdependence, and rules about missing data. Small samples will produce imprecise estimates. If assignment is impossible, carefully document comparison limitations and avoid presenting the result as randomized evidence. The institution need not solve all of educational science before launching; it should steadily increase the credibility of what it learns.

\subsection*{70. College Partnership: The Risk of Borrowed Credibility}

Suppose a college invites Bayes to deliver an applied AI elective. A positive faculty meeting and encouraging student feedback do not yet establish quality or a sustainable partnership. The pilot contract should specify which institution authorizes the program, who owns assessment standards, who grades, how student data may be used, and which claims may be published.

Collect baseline tasks for a representative learner population; measure completion and independent competence; document how college faculty participate and whether their teaching capability improves; inspect operational issues; and obtain a formal buyer evaluation after delivery. If the college selects only its strongest students for the pilot, that should be disclosed before generalizing outcomes to the whole campus. If Bayes supplies detailed curriculum but college faculty perform the teaching, the results reflect the joint system; avoid attributing everything to 1 partner alone.

\subsection*{71. Enterprise Engagement: The Real Unit Is the Workflow}

Suppose a company wants analysts to use AI more competently. Define 1 approved workflow: handling a standard report, evaluating a knowledge retrieval process, or checking the quality of an internal model-assisted response. Collect baseline time, correctness, exception rates, manual review burden, and security constraints. Train a defined group, allow application under legitimate conditions, and measure the workflow again.

If the company changes software, staffing, policies, or input mix at the same time, acknowledge those rival explanations. If workload changes, compare like with like. If employees stop reporting errors because incentives reward faster completion, the metrics may improve while the process becomes worse. Include error-quality and control measures alongside throughput. The customer should understand exactly which observed improvements Bayes can claim and which remain joint outcomes dependent on employer implementation.

\subsection*{72. Measure What Accumulates as Institutional Capital}

A tiny institute should seek more than attractive pilot statistics. It should develop assets that make the next cohort or partnership better: tested tasks, calibrated rubrics, common error taxonomies, mentor-training materials, verified case studies, privacy-respecting evidence practices, and lessons about what learners can sustainably do while employed.

This creates an institutional form of compounding. 1 learner's failure can improve a lesson. 1 college's grading inconsistency can improve the rater guide. 1 enterprise client's implementation bottleneck can sharpen eligibility for the next engagement. The valuable asset is not a growing database for its own sake; it is a \textbf{repeatable ability to make better educational decisions with credible evidence}.

\subsection*{73. What Bayes Should Publish, and When}

At first, publish the teaching method, instructor competence, target capabilities, sample tasks, assessment rubric, limitations, and clearly labeled founder-created demonstration work. After genuine pilots, publish counts and denominators, representative consented artifacts, credible independent task outcomes, and candid accounts of incomplete cases. Do not cherry-pick only top performers while presenting them as typical.

After repeated cohorts with stable assessments, Bayes may responsibly describe trends, again disclosing changes in entry populations and scoring. After stronger comparative studies, it may describe bounded causal findings. Claims about long-run careers, salary changes, research impact, or institutional prestige require additional forms of evidence, often beyond the institute's direct control.

A useful rule for every public number is: \textbf{Could an informed skeptic reproduce the calculation and understand the claim's limits?} If the answer is no, the number is not ready to carry the institution's reputation.
